\begin{textsample}{NEG Dim 2 – 2017\_c5 – Score -3.00 – 2017\_c5/enade\_2017\_qui\_5a144fb539.txt}  \label{ex:f2_neg_050}
CAPÍTULO 5 DISTRIBUIÇÃO DOS CONCEITOS 5.1 PANORAMA NACIONAL DA DISTRIBUIÇÃO DOS CONCEITOS O Conceito Enade é um indicador de qualidade que avalia os cursos por intermédio dos desempenhos dos estudantes a partir dos resultados do Enade.
Os valores possíveis vão de 1 ( pior situação ) a 5 ( melhor situação ).
Os cursos com apenas um ou sem nenhum concluinte participante não obtêm o Conceito Enade, ficando Sem Conceito ( SC ).
A Tabela 5.1a apresenta a quantidade e distribuição de cursos de Licenciatura de Química participantes do Enade/2017, por faixa de conceito e Grande Região. \textbf{Enfatiza} -se, mais uma vez, que a diferença entre os cursos tabulados neste capítulo e no capítulo 2 corresponde aos cursos Sem Conceito, ou seja àqueles, em princípio, sem alunos concluintes que houvessem participado da prova ou àqueles que tivessem tão somente um aluno concluinte¹⁸.
Observando -se os dados da Tabela 5.1a, nota -se que, dos 263 cursos participantes, 89 ( 33,8 \% ) classificaram -se com conceito 3, o valor modal.
Esse foi também o conceito modal nas regiões : Centro-Oeste com 43,8 \% e Sudeste com 32,6 \% ( nesta região o conceito 4 também obteve o mesmo percentual que o 3 ).
Nas regiões Norte e Nordeste, a moda foi o conceito 2, atribuído para 56,5 \% e 41,0 \% dos cursos da região, respectivamente.
E nas regiões Sul e Sudeste o conceito 4 foi o modal com 40,9 \% e 32,6 \%.
O conceito 2 foi o segundo mais frequente em nível nacional ( 27,0 \%, correspondendo a 71 cursos ) e o conceito 4, o terceiro ( 25,1 \%, correspondendo a 66 cursos ).
Houve, ainda, 16 cursos ( 6,1 \% ) que receberam conceito 5 e 11 cursos ( 4,2 \% ) que receberam conceito 1.
Dos 263 cursos de Química ( Licenciatura ), dez ( 3,8 \% ) ficaram Sem Conceito ( SC ). 18 Estes cursos com somente um concluinte não foram considerados no capítulo 2.
A região Norte participou com 23 cursos ou 8,7 \% do total nacional.
Desses, 13 cursos receberam o conceito 2, o que equivale a 56,5 \% do total regional, além de ser o conceito modal na região, como já comentado.
A quatro cursos ( 17,4 \% ) atribuiu -se conceito 1 e a cinco cursos ( 21,7 \% ), conceito 3.
Um curso ficou Sem Conceito ( SC ), e nenhum curso recebeu conceito 4 ou 5 nessa região.
A região Nordeste participou com 78 cursos ou 29,7 \% do total nacional.
Nessa região, 32 cursos ( 41,0 \% em termos regionais ) obtiveram conceito 2, o conceito modal para a região, como já comentado.
Os conceitos 1 e 3 foram atribuídos, respectivamente, a três e a 27 cursos ( 3,8 \% e 34,6 \%, respectivamente ).
Ainda nessa região, 14 cursos ( 17,9 \% ) foram avaliados com o conceito 4, e um curso ( 1,3 \% ), com conceito 5.
Um curso ficou Sem Conceito ( SC ).
Dos 86 cursos participantes da região Sudeste ( 32,7 \% do total nacional ), os conceitos 3 e 4 foram os conceitos modais, como já anteriormente assinalado, tendo sido obtidos, por 28 cursos cada ( 32,6 \% ).
O conceito 1 foi atribuído a três cursos ( 3,5 \% ), e o conceito 2, a 12 cursos ( 14,0 \% ).
Receberam o conceito 5 dez cursos ( 11,6 \% ).
Nessa região, cinco cursos ficaram Sem Conceito ( SC ).
Os 44 cursos da região Sul corresponderam a 16,7 \% do total nacional.
Como já indicado, houve predominância do conceito 4, atribuído a 40,9 \% dos cursos da região ( conceito modal ), ou seja, atribuído a 18 dos 44 cursos participantes na região Sul.
O conceito 2 foi atribuído a cinco cursos ( 11,4 \% ), o conceito 3, a 15 cursos ( 34,1 \% ) e o conceito 5, a cinco cursos ( 11,4 \% ).
Um curso ficou Sem Conceito ( SC ) na região Sul, e nenhum recebeu conceito 1.
Já dos 32 cursos participantes na região Centro-Oeste ( 12,2 \% do total nacional ), 14 receberam conceito 3, o conceito modal ( 43,8 \% ), como já fora destacado.
A um curso ( 3,1 \% ) foi atribuído conceito 1, a nove cursos ( 28,1 \% ), conceito 2, e a seis cursos ( 18,8 \% ), conceito 4.
Nessa região, dois cursos ficaram Sem Conceito ( SC ) e nenhum recebeu conceito 5.
O Gráfico 5.1a apresenta a distribuição cumulativa do Conceito Enade dos cursos de Licenciatura segundo a Grande Região.
Num gráfico de distribuição cumulativa, em cada valor de abscissa ( eixo x ), o valor apresentado no eixo das ordenadas ( eixo y ) é o valor acumulado da variável.
Por exemplo, a região Norte ( linha azul ) apresenta 100 \% dos cursos com conceito 3 ou abaixo deste valor, pois a poligonal cumulativa alcança 100 \% no conceito 3.
A região Sul ( linha roxa ), por outro lado, apresenta um pouco mais de 46 \% dos cursos na mesma situação : com conceito 3 ou abaixo.
Na região Centro-Oeste ( linha preta ), todos os cursos têm conceitos 4 ou abaixo disso, já que a poligonal cumulativa alcança 100 \% para o conceito 4.
Em linhas gerais, em um gráfico de distribuição cumulativa, poligonais mais à esquerda representam Grandes Regiões com pior distribuição, e poligonais mais à direita, Grandes Regiões com melhores conceitos.
A região Norte ( linha azul ) apresenta os cursos com conceitos mais baixos, ao passo que as regiões Sul ( linha roxa ) e Sudeste ( linha vermelha ) apresentam as distribuições com valores maiores.
Cursos em IES situadas nas regiões Nordeste ( linha verde ) e Centro-Oeste ( linha preta ) estão em uma situação intermediária. melhores conceitos.
A região Norte ( linha azul ) apresenta os cursos com conceitos mais baixos, ao passo que as regiões Sul ( linha roxa ) e Sudeste ( linha vermelha ) apresentam as distribuições com valores maiores.
Cursos em IES situadas nas regiões Nordeste ( linha verde ) e Centro-Oeste ( linha preta ) estão em uma situação intermediária.
A Tabela 5.1b apresenta a quantidade e distribuição de cursos de Bacharelado de Química participantes do Enade/2017, por faixa de conceito e Grande Região. \textbf{Enfatiza} -se, mais uma vez, que a diferença entre os cursos tabulados neste capítulo e no capítulo 2 corresponde aos cursos sem conceito, ou seja àqueles, em princípio, sem alunos concluintes que houvessem participado da prova ou àqueles que tivessem tão somente um aluno concluinte19.
Observando -se os dados da Tabela 5.1b, nota -se que, dos 128 cursos participantes, 51 ( 39,8 \% ) classificaram -se com conceito 3, o valor modal.
Esse foi também o conceito modal em quase todas as regiões : Nordeste ( 41,2 \% ), Sudeste ( 33,9 \% ), Sul ( 46,9 \% ) e Centro-Oeste ( 53,8 \% ).
Na região Norte, a moda foi o conceito 1, atribuído a 50,0 \% dos cursos da região.
O 19 Estes cursos com somente um concluinte não foram considerados no capítulo 2. conceito 2 foi o segundo mais frequente em nível nacional ( 28,9 \%, correspondendo a 37 cursos ) e o conceito 4, o terceiro ( 14,1 \%, correspondendo a 18 cursos ).
Houve, ainda, 11 cursos ( 8,6 \% ) que receberam conceito 1, e oito cursos ( 6,3 \% ) que receberam conceito 1.
Dos 128 cursos de Química ( Bacharelado ), três ( 2,3 \% ) ficaram Sem Conceito ( SC ).
A região Norte participou com quatro cursos ou 3,1 \% do total nacional.
Desses, dois cursos receberam o conceito 1, o que equivale a 50,0 \% do total regional, além de ser o conceito modal na região, como já comentado.
A um curso ( 25,0 \% ) atribuiu -se conceito 2, e a outro um curso, o conceito 3.
Nenhum curso recebeu os demais conceitos nessa região ou ficou Sem Conceito ( SC ).
A região Nordeste participou com 17 cursos ou 13,3 \% do total nacional.
Nessa região, sete cursos ( 41,2 \% em termos regionais ) obtiveram conceito 3, o conceito modal para a região, como já comentado.
Os conceitos 1 e 2 foram atribuídos, respectivamente, a três e a quatro cursos ( 17,6 \% e 23,5 \%, respectivamente ).
Ainda nessa região, dois cursos ( 11,8 \% ) foram avaliados com o conceito 4.
Um curso ficou Sem Conceito ( SC ), e nenhum recebeu conceito 5.
Dos 62 cursos participantes da região Sudeste ( 48,4 \% do total nacional ), o conceito 3 foi o conceito modal, como já anteriormente assinalado, tendo sido obtido por 21 cursos ( 33,9 \% ).
O conceito 1 foi atribuído a cinco cursos ( 8,1 \% ) e o conceito 2, a 19 cursos ( 30,6 \% ).
Receberam os conceitos 4 e 5, respectivamente, dez cursos ( 16,1 \% ) e seis cursos ( 9,7 \% ).
Nessa região, um curso ficou Sem Conceito ( SC ).
Os 32 cursos da região Sul corresponderam a 25,0 \% do total nacional.
Como já indicado, houve predominância do conceito 3, atribuído a 46,9 \% dos cursos da região ( conceito modal ), ou seja, atribuído a 15 dos 32 cursos participantes na região Sul.
O conceito 1 foi atribuído a um curso ( 3,1 \% ), o conceito 2, a nove cursos ( 28,1 \% ), o conceito 4, a cinco cursos ( 15,6 \% ) e o conceito 5, a um curso ( 3,1 \% ).
Um curso ficou Sem Conceito ( SC ) na região Sul.
Já dos 13 cursos participantes na região Centro-Oeste ( 10,2 \% do total nacional ), sete receberam conceito 3, o conceito modal ( 53,8 \% ), como já fora destacado.
A quatro cursos ( 30,8 \% ) foi atribuído conceito 2, a um curso ( 7,7 \% ), conceito 4, e a um curso ( 7,7 \% ), o conceito 5.
Nessa região, nenhum curso ficou Sem Conceito ( SC ) ou recebeu conceito 1.
O Gráfico 5.1b apresenta a distribuição cumulativa do Conceito Enade dos cursos de Bacharelado segundo a Grande Região.
Num gráfico de distribuição cumulativa, em cada valor de abscissa ( eixo x ), o valor apresentado no eixo das ordenadas ( eixo y ) é o valor acumulado da variável.
Por exemplo, a região Norte ( linha azul ) apresenta 100 \% dos cursos com conceito 3 ou abaixo deste valor, já que a poligonal cumulativa alcança 100 \% no conceito 3.
A região Sudeste ( linha vermelha ), por outro lado, apresenta um pouco mais de 73 \% dos cursos na mesma situação : com conceito 3 ou abaixo.
Na região Nordeste ( linha verde ), todos os cursos têm conceitos 4 ou abaixo disso, já que a poligonal cumulativa alcança 100 \% para o conceito 4.
Em linhas gerais, num gráfico de distribuição cumulativa, poligonais mais à esquerda representam Grandes Regiões com pior distribuição, e poligonais mais à direita, Grandes Regiões com melhores conceitos.
A região Norte ( linha azul ) apresenta os cursos com conceitos mais baixos, ao passo que as regiões Sudeste e Centro-Oeste ( linha preta ) apresentam as distribuições com valores maiores, mas não é possível observar uma hierarquia. 5.2 CONCEITOS POR CATEGORIA ADMINISTRATIVA, POR MODALIDADE DE ENSINO E GRANDE REGIÃO A Tabela 5.2a apresenta a distribuição dos cursos de Licenciatura de Química participantes do Enade/2017, por Categoria Administrativa e Modalidade de ensino, de acordo com os conceitos por eles alcançados, segundo as Grandes Regiões brasileiras.
Dos 263 cursos participantes, 214 ( 81,4 \% ) eram ministrados em Instituições Públicas e 49 ( 18,6 \% ), em Instituições Privadas.
Quanto à Modalidade de ensino, 92,8 \% dos cursos eram Presenciais e 7,2 \% a Distância.
De acordo com as informações da Tabela 5.2a, em termos nacionais, entre Instituições Públicas e Privadas, observa -se que, dos 16 cursos avaliados com conceito 5, todos eram oferecidos em IES Públicas.
Dos 214 cursos participantes de IES Públicas, o conceito 3 foi o valor modal, atribuído a 70 cursos ( 32,7 \% ).
Entre os demais cursos de Instituições Públicas participantes, sete cursos ( 3,3 \% ) foram avaliados com conceito 1, 57 cursos ( 26,6 \% ) obtiveram conceito 2, 58 cursos ( 27,1 \% ) receberam conceito 4 e, como já comentado, 16 cursos ( 7,5 \% ) receberam conceito 5.
Nesta categoria, seis cursos ( 2,8 \% ) ficaram Sem Conceito ( SC ).
Na rede privada, o conceito modal também foi 3, atribuído a 19 cursos, o correspondente a 38,8 \% dos 49 cursos da categoria.
Entre os demais cursos participantes, quatro ( 8,2 \% ) receberam conceito 1, e 14 ( 28,6 \% ), conceito 2.
O conceito 4 foi atribuído a oito cursos ( 16,3 \% ), e o conceito 5, a nenhum curso.
Nesta Categoria Administrativa, quatro cursos ( 8,2 \% ) ficaram Sem Conceito ( SC ).
Dos 92,8 \% cursos participantes da modalidade presencial, o conceito 3 foi o valor modal, atribuído a 33,2 \% dos cursos.
Entre os demais cursos nesta modalidade de ensino, 25,4 \% dos cursos foram avaliados com conceito 4, 27,0 \% obtiveram conceito 2, 6,1 \% receberam conceito 5 e 4,5 \% receberam conceito 1.
Nesta modalidade 3,7 \% dos cursos ficaram Sem Conceito ( SC ).
Entre os cursos a Distância, o conceito modal foi também 3, atribuído a 42,1 \% dos cursos.
Entre os demais cursos participantes, 26,3 \% receberam conceito 2, 21,1 \%, conceito 4 e 5,3 \%, conceito 5.
Nessa Modalidade 5,3 \% dos cursos ficaram Sem Conceito ( SC ), e nenhum recebeu conceito 1.
Na análise por região, observa -se que, na região Norte, as Instituições Públicas participaram com 19 cursos ( 82,6 \% do total regional ), dos quais um ficou Sem Conceito ( SC ), e nenhum recebeu conceito 4 ou 5.
O conceito modal para as Instituições Públicas na região foi 2, atribuído a dez cursos, correspondendo a 52,6 \% dessa categoria na região.
Cinco cursos foram avaliados com conceito 3 ( 26,3 \% ), e três cursos, com conceito 1.
As Instituições Privadas participaram com quatro cursos na região Norte ( 17,4 \% do total regional ), dos quais três obtiveram conceito 2, o valor modal, e o conceito 1 foi atribuído a um curso.
Nenhum curso ficou Sem Conceito ( SC ) ou com conceito 3 ou 4 ou 5 nessa combinação de categoria e região.
A maioria dos cursos oferecidos na região Norte eram presenciais, 91,3 \%, e eles foram \textbf{distribuídos} da seguinte forma : 52,4 \% dos cursos ficaram com conceito 2 ( conceito modal ), 23,8 \% ficaram com conceito 3 e 19,0 \%, com conceito 1.
Nota -se ainda que 4,8 \% ficaram Sem Conceito ( SC ).
Dos cursos oferecidos na modalidade a distância, todos receberam o conceito 2, e nenhum ficou Sem Conceito ( SC ).
Na região Nordeste, a rede pública concentrou 74 dos 78 cursos participantes, o equivalente a 94,9 \% do total da região.
Dentre os cursos oferecidos por IES Públicas no Nordeste, 30 cursos foram avaliados com conceito 2, o conceito modal.
Três cursos obtiveram conceito 1, 26 cursos receberam conceito 3, 13 cursos ficaram com conceito 4, e um curso alcançou o conceito máximo, 5.
Nessa combinação de categoria e região, um curso ficou Sem Conceito ( SC ).
As Instituições Privadas da região Nordeste participaram com quatro cursos ( 5,1 \% do total da região ).
Desses, dois obtiveram conceito 2, o valor modal, seguidos por outros dois cursos, que receberam conceito 3 e 4 ( um curso cada ).
Nenhum curso ficou Sem Conceito ( SC ) ou com conceito 1 ou 5 nessa região.
A maioria dos cursos na região Nordeste foram presencias : 93,6 \%.
Para estes, a moda foi o conceito 2 com 42,5 \% dos cursos.
Os demais cursos presenciais da região foram alocados aos conceitos 1 ( 4,1 \% ), 3 ( 34,2 \% ), 4 ( 16,4 \% ) e 5 ( 1,4 \% ).
Dos cursos presenciais, 1,4 \% ficou Sem Conceito ( SC ).
Os cursos a Distância receberam o conceito 3 ( dois cursos ), 4 ( dois cursos ) e 2 ( um curso ).
Nenhum curso a Distância da região ficou Sem Conceito ( SC ) ou recebeu os conceitos 1 e 5.
Na região Sudeste, a proporção de cursos da rede privada, 31,4 \%, foi mais \textbf{elevada} do que a registrada nas demais regiões brasileiras, correspondendo a 27 dos 86 cursos participantes.
Nessa categoria e região, o conceito modal foi 3, atribuído a 11 cursos.
Dois cursos ficaram Sem Conceito ( SC ).
Os demais foram avaliados com conceito 1 ( três cursos ), conceito 2 ( oito cursos ), conceito 4 ( três cursos ) e conceito 5 ( nenhum curso ).
Entre os 59 cursos oferecidos por Instituições Públicas na região Sudeste, o conceito modal foi 4, atribuído a 25 cursos.
Os demais cursos receberam os conceitos 2 ( quatro cursos ), 3 ( 17 cursos ) e 5 ( dez cursos ).
Na região Sudeste, nesta categoria, nenhum curso recebeu o conceito 1, e três cursos ficaram Sem Conceito ( SC ).
Também na região Sudeste, a maioria dos cursos foi presencial : 87,2 \%.
Para estes, a moda foi o conceito 4 com 34,7 \% dos cursos.
Os demais cursos presenciais da região foram alocados aos conceitos 1 ( 4,0 \% dos cursos ), 2 ( 13,3 \% ), 3 ( 29,3 \% ) e 5 ( 12,0 \% ).
Cinco cursos presenciais ficaram Sem Conceito ( SC ).
Para os cursos a Distância, a moda foi o conceito 3 ( 54,5 \% ).
Os demais cursos a Distância da região foram alocados aos conceitos 2 ( 18,2 \% dos cursos ), 4 ( 18,2 \% ) e 5 ( 9,1 \% ).
Nenhum curso a Distância da região foi alocado ao conceito 1 ou ficou Sem Conceito ( SC ).
As Instituições Públicas concentraram 34 dos 44 cursos participantes da região Sul, 77,3 \% do total regional.
Desses, 14 obtiveram conceito 4, o conceito modal.
Nesta combinação de Categoria Administrativa e Grande Região, os conceitos 2 e 3 foram atribuídos a cinco e a dez cursos, respectivamente.
Cinco cursos receberam conceito 5, nenhum recebeu conceito 1 ou ficou Sem Conceito ( SC ).
As Instituições Privadas na região Sul participaram com dez cursos, 22,7 \% dos cursos da região.
O conceito modal foi 3, atribuído a cinco cursos.
Aos demais cursos foram atribuídos o conceito 4 ( quatro cursos ).
Nenhum curso da região Sul, oferecido por IES Privadas, recebeu os demais conceitos e um curso ficou Sem Conceito ( SC ).
Na região Sul, todos os cursos foram oferecidos na modalidade presencial.
Na região Centro-Oeste, 28 dos 32 cursos participantes eram de Instituições Públicas ( 87,5 \% em termos regionais ).
Desses, 12 receberam conceito 3, o conceito modal.
Entre os demais, um recebeu conceito 1, oito receberam conceito 2 e seis receberam conceito 4.
Nenhum curso ficou com conceito 5, e um curso ficou Sem Conceito ( SC ).
Quanto aos quatro cursos oferecidos por Instituições Privadas na região Centro-Oeste ( 12,5 \% do total regional ), dois cursos apresentaram conceito 3, conceito modal, e um ficou com conceito 2.
Nenhum curso ficou com os demais conceitos, e um ficou Sem Conceito ( SC ).
Nesta região, 96,9 \% dos cursos foram presenciais e a moda foi o conceito 3 com 45,2 \% dos cursos.
Os demais foram alocados aos conceitos 1 ( 3,2 \% ), 2 ( 29,0 \% ) e 4 ( 19,4 \% ).
Nenhum curso presencial ficou com conceito 5, e um curso ficou Sem Conceito ( SC ).
O único curso da modalidade a distância ( 3,1 \% ) ficou Sem Conceito ( SC ).
O Gráfico 5.2a apresenta a distribuição cumulativa do Conceito Enade dos cursos de Licenciatura segundo a Categoria Administrativa da IES.
Os cursos em IES Públicas ( linha azul ) apresentam uma poligonal mais à direita do que os em IES Privadas ( linha verde ) e, consequentemente, uma distribuição de Conceitos Enade com valores maiores.
Corroborando essa informação, nota -se que todos os cursos em IES Privadas apresentam conceito 4 ou abaixo disso, já que a poligonal cumulativa alcança 100 \% no conceito 4.
A Tabela 5.2b apresenta a distribuição dos cursos de Bacharelado participantes do Enade/2017 de Química, por Categoria Administrativa e Modalidade de ensino, de acordo com os conceitos por eles alcançados, segundo as Grandes Regiões brasileiras.
Dos 128 cursos participantes, 86 ( 67,2 \% ) eram ministrados em Instituições Públicas e 42 ( 32,8 \% ), em Instituições Privadas.
Quanto à Modalidade de ensino, todos os cursos eram Presenciais.
De acordo com as informações da Tabela 5.2b, em termos nacionais, entre Instituições Públicas e Privadas, observa -se que, dos oito cursos avaliados com conceito 5, sete eram oferecidos em IES Públicas e um, em IES Privadas.
Dos 86 cursos participantes de IES Públicas, o conceito 3 foi o valor modal, atribuído a 41 cursos ( 47,7 \% ).
Entre os demais cursos de Instituições Públicas participantes, seis cursos ( 7,0 \% ) foram avaliados com conceito 1, 13 cursos ( 15,1 \% ) obtiveram conceito 2, 18 cursos ( 20,9 \% ) receberam conceito 4 e, como já comentado, sete cursos ( 8,1 \% ) receberam conceito 5.
Nesta categoria, um curso ( 1,2 \% ) ficou Sem Conceito ( SC ).
Na rede privada, o conceito modal foi 2, atribuído a 24 cursos, o correspondente a 57,1 \% dos 42 cursos da categoria.
Entre os demais cursos participantes, cinco ( 11,9 \% ) receberam conceito 1, e dez ( 23,8 \% ), conceito 3.
O conceito 5 foi atribuído a um curso ( 2,4 \% ) e o conceito 4 não foi atribuído a nenhum curso.
Nesta Categoria Administrativa, dois cursos ( 4,8 \% ) ficaram Sem Conceito ( SC ).
Como já informado, todos os cursos de Química ( Bacharelado ) foram vinculados á modalidade presencial, logo a distribuição conceitual é igual à descrita ao nível nacional.
Na análise por região, observa -se que, na região Norte, as Instituições Públicas participaram com todos os quatro cursos, dos quais nenhum recebeu conceito 4 ou 5.
O conceito modal para as Instituições Públicas na região foi 1, atribuído a dois cursos, correspondendo a 50,0 \% dessa categoria na região, e dois cursos foram avaliados com conceitos 2 e 3 ( um curso, cada ).
Na região Nordeste, a rede pública também concentrou todos os 17 cursos participantes da região.
Dentre eles, sete cursos foram avaliados com conceito 3, o conceito modal.
Três cursos obtiveram conceito 1, quatro cursos receberam conceito 2 e dois cursos ficaram com conceito 4.
Nessa combinação de categoria e região, um curso ficou Sem Conceito ( SC ), e nenhum curso recebeu o conceito máximo, 5.
Na região Sudeste, a proporção de cursos da rede privada, 43,5 \%, foi mais \textbf{elevada} do que a registrada nas demais regiões brasileiras, correspondendo a 27 dos 62 cursos participantes.
Nessa categoria e região, o conceito modal foi 2, atribuído a 17 cursos.
Um curso ficou Sem Conceito ( SC ), e um ficou com conceito 5.
Os demais foram avaliados com conceito 1 ( quatro cursos ), conceito 3 ( quatro cursos ) e conceito 4 ( nenhum curso ).
Entre os 35 cursos oferecidos por Instituições Públicas na região Sudeste, o conceito modal foi 3, atribuído a 17 cursos.
Os demais cursos receberam os conceitos 1 ( um curso ), 2 ( dois cursos ), 4 ( dez cursos ) e 5 ( cinco cursos ).
Na região Sudeste, nesta categoria, nenhum curso ficou Sem Conceito ( SC ).
As Instituições Públicas concentraram 19 dos 32 cursos participantes da região Sul, 59,4 \% do total regional.
Desses, dez obtiveram conceito 3, o conceito modal.
Nesta combinação de Categoria Administrativa e Grande Região, o conceito 2 foi atribuído a três cursos, o conceito 4, a cinco cursos e o conceito 5, a um curso.
Nenhum curso ficou Sem Conceito ( SC ) ou recebeu conceito 1.
As Instituições Privadas na região Sul participaram com 13 cursos, 40,6 \% dos cursos da região.
O conceito modal foi 2, atribuído a seis cursos.
Aos demais cursos foram atribuídos os conceitos 1 ( um curso ) e 3 ( cinco cursos ).
Nenhum curso da região Sul oferecido por IES Privadas recebeu conceitos 4 ou 5, e um curso ficou Sem Conceito ( SC ).
Na região Centro-Oeste, 11 dos 13 cursos participantes eram de Instituições Públicas ( 84,6 \% em termos regionais ).
Desses, seis receberam conceito 3, o conceito modal.
Entre os demais, três receberam conceito 2, um recebeu conceito 4, e um recebeu conceito 5.
Nenhum curso ficou Sem Conceito ( SC ) ou com conceito 1.
Quanto aos dois cursos oferecidos por Instituições Privadas na região Centro-Oeste ( 15,4 \% do total regional ), os conceitos 2 e 3 foram atribuídos a um curso cada.
Nenhum curso ficou Sem Conceito ( SC ) ou com conceito 1, 4 ou 5.
O Gráfico 5.2b apresenta a distribuição cumulativa do Conceito Enade dos cursos de Bacharelado segundo a Categoria Administrativa da IES.
Os cursos em IES Públicas ( linha azul ) apresentam uma poligonal mais à direita do que os em IES Privadas ( linha verde ) e, consequentemente, uma distribuição de Conceitos Enade com valores maiores.
Corroborando essa informação, nota -se que os cursos de IES Privadas não receberam conceito 4, já que a poligonal cumulativa desses cursos se configura como uma reta horizontal entre os conceitos 3 e 4, ou seja, não houve acréscimo do percentual de cursos quando se analisa o conceito 4, o que é totalmente diferente, quando se compara a poligonal dos cursos em IES Públicas, que obteve um acréscimo de 21,2 \% quando analisados os cursos com conceito 4. 5.3 CONCEITOS POR ORGANIZAÇÃO ACADÊMICA E GRANDE REGIÃO Na Tabela 5.3a, encontra -se a distribuição dos conceitos atribuídos aos cursos de Licenciatura participantes do Enade/2017 na Área de Química, por Organização Acadêmica, segundo as Grandes Regiões brasileiras.
Dos 263 cursos participantes, 170 eram oferecidos em Universidades, 71, em CEFET/IFET, 13, em Faculdades e nove, em Centros Universitários.
Esta distribuição corresponde a, respectivamente, 64,6 \%, 27,0 \%, 4,9 \% e 3,4 \% dos cursos.
De acordo com os dados apresentados, de todos os 16 cursos avaliados com conceito 5, 15 eram vinculados a Universidades.
Os cursos vinculados a esse tipo de Organização Acadêmica tiveram o conceito 3 como conceito modal, atribuído a 57 cursos ( 33,5 \% ).
Os demais cursos vinculados a Universidades receberam os conceitos 1 ( cinco cursos ), 2 ( 34 cursos ), 4 ( 51 cursos ) e conceito 5 ( 15 cursos, como já mencionado ).
Ficaram Sem Conceito ( SC ) oito cursos.
Entre os cursos em Centros Universitários, o conceito modal foi 2, atribuído a cinco cursos ( 55,6 \% ).
Nenhum curso vinculado a esse tipo de Organização Acadêmica ficou Sem Conceito ( SC ) ou recebeu conceito 4 ou 5, e os outros cursos receberam os conceitos 1 ( dois cursos ) e 3 ( dois cursos ).
Dos 13 cursos mantidos por Faculdades, cinco ( 38,5 \% ) receberam conceito 2, o conceito modal.
O segundo conceito mais frequentemente atribuído aos cursos oferecidos por essa categoria de IES foi o conceito 3, recebido por quatro cursos.
Na sequência, dois cursos obtiveram conceito 4, e um curso, conceito 1.
Nenhum curso foi atribuído ao conceito 5, e um curso ficou Sem Conceito ( SC ).
Dos 71 cursos mantidos por CEFET/IFET, 27 ( 38,0 \% ) receberam conceito 2, o conceito modal.
O segundo conceito mais frequentemente atribuído aos cursos oferecidos por essa categoria de IES foi o conceito 3, recebido por 26 cursos.
Foram também 13 os cursos que receberam conceito 4, três, o conceito 1 e um, o conceito 5.
Um curso ficou Sem Conceito ( SC ).
Considerando -se separadamente as regiões brasileiras, verifica -se que, na região Norte, as Universidades participaram com 13 dos 23 cursos participantes.
Sete cursos receberam conceito 2, conceito modal.
Os demais cursos receberam os conceitos 1 ( dois cursos ) e 3 ( três cursos ).
Nenhum curso ficou com conceito 4 ou 5, e um curso ficou Sem Conceito ( SC ).
Os Centros Universitários da região Norte foram representados por um curso, ao qual foi atribuído o conceito 2.
As Faculdades participaram com três cursos na região Norte que receberam os conceitos 2 ( dois cursos, conceito modal ) e 1 ( um curso ).
Nenhuma das Faculdades da região ficou Sem Conceito ( SC ) ou obteve conceito 3, 4 ou 5.
Os CEFET/IFET participaram com seis cursos, sendo três com conceito 2 ( conceito modal ), dois com conceito 3 e um com conceito 1.
Nenhum curso ficou Sem Conceito ( SC ) ou recebeu conceito 4 ou 5, nesta organização acadêmica, desta região.
Na região Nordeste, as Universidades participaram com 44 dos 78 cursos da Área de Química ( Licenciatura ) oferecidos na região, dos quais um ficou Sem Conceito ( SC ).
Nota -se que tais cursos apresentaram o conceito 3 como modal, atribuído a 16 cursos.
Obtiveram os conceitos 1, 2, 4 e 5, respectivamente, dois, 15, nove e um cursos.
Os Centros Universitários contaram com um curso, que foi conceituado com 3.
As Faculdades foram representadas por três cursos na região Nordeste, dois receberam conceito 2, e um recebeu conceito 4.
Na região Nordeste, nenhum curso oferecido por Faculdades ficou Sem Conceito ( SC ).
Os CEFET/IFET participaram com 30 cursos, dos quais 15 receberam conceito 2 ( conceito modal ), dez receberam conceito 3, quatro receberam conceito 4 e um recebeu conceito 1.
Nenhum curso ficou Sem Conceito ( SC ), nem recebeu conceito 5.
Na região Sudeste, as Universidades concentraram 61 dos 86 cursos de Química ( Licenciatura ) da região.
Entre os cursos em Universidades, nessa região, o conceito modal foi 4, atribuído a 22 cursos.
Os demais cursos receberam os conceitos 1 ( um curso ), 2 ( 6 cursos ), 3 ( 18 cursos ) e 5 ( nove cursos ).
Cinco cursos ficaram Sem Conceito ( SC ).
Os Centros Universitários participaram com seis cursos na região Sudeste, dos quais quatro obtiveram o conceito modal 2, e dois cursos obtiveram conceito 1.
Nenhum curso ficou Sem Conceito ( SC ) ou recebeu conceito 3, 4 ou 5 nessa combinação de organização acadêmica e região.
As Faculdades também foram representadas por seis cursos na região Sudeste, dos quais nenhum ficou Sem Conceito ( SC ) ou recebeu conceito 1 ou 5.
Aos demais cursos foram atribuídos os conceitos 2 ( um curso ), 3 ( quatro cursos, o conceito modal ) e 4 ( um curso ).
Os CEFET/IFET participaram com 13 cursos, os quais receberam conceitos 3 ( seis cursos, o conceito modal ), 4 ( cinco cursos ), 2 ( um curso ) e 5 ( um curso ).
Nenhum curso ficou Sem Conceito ( SC ), nem recebeu conceito 1.
Dos 44 cursos da região Sul, 36 eram de Universidades, e o conceito modal 4 foi atribuído a 17 cursos.
Nesse tipo de organização, um curso ficou Sem Conceito ( SC ), e nenhum curso recebeu conceito 1.
Os demais cursos receberam os conceitos 2 ( um curso ), 3 ( 12 cursos ) e 5 ( cinco cursos ).
Nenhum curso representou os Centros Universitários e as Faculdades da região Sul.
Os CEFET/IFET foram representados por oito cursos na região, dos quais quatro receberam conceito 2 ( conceito modal ), três receberam conceito 3 e um recebeu conceito 4.
Nenhum curso ficou Sem Conceito ( SC ) ou recebeu conceito 1 ou 5, nesta organização e região.
Na região Centro-Oeste, 16 dos 32 cursos eram de Universidades.
Nesse tipo de organização, nenhum curso ficou Sem Conceito ( SC ) ou recebeu conceito 1 ou 5.
O conceito 2 foi atribuído a cinco cursos, o conceito 3, a oito cursos ( conceito modal ) e o conceito 4, a três cursos.
Os Centros Universitários da região Centro-Oeste contaram com um curso, que obteve o conceito 3.
O único curso representando as Faculdades na região Centro-Oeste ficou Sem Conceito ( SC ).
Dos 14 cursos representando os CEFET/IFET, cinco receberam conceito 3 ( conceito modal ), um recebeu conceito 1, quatro, o conceito 2 e três, o conceito 4.
Apenas um curso em CEFET/IFET ficou Sem Conceito ( SC ) e a nenhum curso foi atribuído o conceito 5.
O Gráfico 5.3a apresenta a distribuição cumulativa do Conceito Enade dos cursos de Licenciatura, segundo a Organização Acadêmica da IES.
Os cursos em Universidades ( linha azul ) apresentam uma poligonal mais à direita do que as demais, denotando uma melhor distribuição de Conceitos Enade ( 59,3 \% dos cursos obtiveram conceito 3 ou abaixo deste valor ).
No outro extremo, temos os cursos oferecidos em Centros Universitários ( linha verde ) com a poligonal mais à esquerda e, consequentemente, uma distribuição de Conceitos Enade com menores valores ( 100 \% dos cursos obtiveram conceito 3 ou abaixo deste valor ).
IES de Faculdades ( linha vermelha ) e CEFET/IFET ( linha roxa ) apresentam uma situação intermediária.
Na Tabela 5.3b, encontra -se a distribuição dos conceitos atribuídos aos cursos de Bacharelado participantes do Enade/2017 na Área de Química, por Organização Acadêmica, segundo as Grandes Regiões brasileiras.
Dos 128 cursos participantes, 110 eram oferecidos em Universidades, oito em Centros Universitários, sete em Faculdades e três em CEFET/IFET.
Esta distribuição corresponde a, respectivamente, 85,9 \%, 6,3 \%, 5,5 \% e 2,3 \% dos cursos.
De acordo com os dados apresentados, todos os oito cursos avaliados com conceito 5 eram vinculados a Universidades.
Os cursos vinculados a esse tipo de Organização Acadêmica tiveram o conceito 3 como conceito modal, atribuído a 45 cursos ( 40,9 \% ).
Os demais cursos vinculados a Universidades receberam os conceitos 1 ( nove cursos ), 2 ( 27 cursos ), 4 ( 18 cursos ) e conceito 5 ( oito cursos, como já mencionado ).
Ficaram Sem Conceito ( SC ) três cursos.
Entre os cursos em Centros Universitários, o conceito modal foi 2, atribuído a cinco cursos ( 62,5 \% ).
Nenhum curso vinculado a esse tipo de Organização Acadêmica ficou Sem Conceito ( SC ) ou recebeu conceito 4 ou 5, e os outros cursos receberam os conceitos 1 ( um curso ) e 3 ( dois cursos ).
Dos sete cursos mantidos por Faculdades, cinco ( 71,4 \% ) receberam conceito 2, o conceito modal.
Outros dois cursos foram conceituados igualmente entre os conceitos 1 e 3.
Os conceitos 4 ou 5 não foram atribuídos a nenhum curso e nenhum curso ficou Sem Conceito ( SC ).
Dos três cursos mantidos por CEFET/IFET, todos receberam conceito 3, ou seja, nenhum curso ficou Sem Conceito ( SC ) ou recebeu os demais conceitos.
Considerando -se separadamente as regiões brasileiras, verifica -se que, na região Norte, as Universidades participaram com todos os quatro cursos participantes na região.
Dois cursos receberam conceito 1, conceito modal.
Os demais cursos receberam os conceitos 2 ( um curso ) e 3 ( um curso ).
Nenhum curso ficou Sem Conceito ( SC ) ou recebeu conceito 4 ou 5 nesta combinação de categoria e região.
Na região Nordeste, as Universidades também participaram com todos os 17 cursos da Área de Química ( Bacharelado ) oferecidos, dos quais um ficou Sem Conceito ( SC ), e nenhum recebeu conceito 5.
Obtiveram os conceitos 1 ( três cursos ), 2 ( quatro cursos ), 3 ( sete cursos, conceito modal ) e 4 ( dois cursos ).
Na região Sudeste, as Universidades concentraram 47 dos 62 cursos de Química ( Bacharelado ).
Entre os cursos em Universidades, nessa região, o conceito modal foi 3, atribuído a 17 cursos.
Os demais cursos receberam os conceitos 1 ( três cursos ), 2 ( dez cursos ), 4 ( dez cursos ) e 5 ( seis cursos ).
Um curso ficou Sem Conceito ( SC ).
Os Centros Universitários participaram com sete cursos na região Sudeste, dos quais cinco obtiveram o conceito modal 2.
O conceito 1 foi atribuído a um curso e o conceito 3, a um curso.
Nenhum curso ficou Sem Conceito ( SC ) ou ficou com conceito 4 ou 5 nessa combinação de organização acadêmica e região.
Já as Faculdades foram representadas por seis cursos na região Sudeste, dos quais nenhum ficou Sem Conceito ( SC ) ou recebeu conceito 4 ou 5.
Aos demais cursos foram atribuídos os conceitos 1 ( um curso ), 2 ( quatro cursos, o conceito modal ) e 3 ( um curso ).
Dois cursos em CEFET/IFET participaram nesta região, os quais receberam conceito 3.
Dos 32 cursos da região Sul, todos eram de Universidades, para os quais o conceito modal foi 3, atribuído a 15 cursos.
Nesse tipo de organização, um curso ficou Sem Conceito ( SC ).
Os demais cursos receberam os conceitos 1 e 5 ( um curso, cada ), 2 ( nove cursos ) e 4 ( cinco cursos ).
Na região Centro-Oeste, dez dos 13 cursos eram de Universidades.
Nesse tipo de organização, nenhum curso ficou Sem Conceito ( SC ) ou recebeu o conceito 1.
Os conceitos 4 e 5 foram atribuídos a um curso cada, o conceito 2, a três cursos e o conceito 3, conceito modal, a cinco cursos.
Os Centros Universitários da região Centro-Oeste, assim como as Faculdades e os CEFET/IFET foram representados com um curso cada, sendo que eles receberam conceitos 3, 2 e 3, respectivamente.
O Gráfico 5.3b apresenta a distribuição cumulativa do Conceito Enade dos cursos de Bacharelado segundo a Organização Acadêmica da IES.
Os cursos em Universidades ( linha azul ) apresentam uma poligonal mais à direita do que as demais, denotando uma melhor distribuição de Conceitos Enade ( 75,7 \% dos cursos receberam conceito 3 ou abaixo deste valor ).
No outro extremo, temos os cursos oferecidos em Faculdades ( linha vermelha ) com a poligonal mais à esquerda e, consequentemente, uma distribuição de Conceitos Enade com menores valores ( 100 \% dos cursos receberam conceito 3 ou abaixo deste valor ).
IES de CEFET/IFET ( linha roxa ) apresentam uma característica bastante singular : todos os cursos receberam conceito 3, pois a poligonal cumulativa parte de 0 \% no conceito 2 e alcança 100 \% no conceito 3.
Todos os cursos oferecidos em Faculdades ( linha vermelha ), Centros Universitários ( linha verde ) e CEFET/IFET ( linha roxa ) receberam conceito 3 ou abaixo.

% matched lemmas: distribuir, elevado, enfatizar
\end{textsample}
